% Appendix B -- Routines written in Pascal.  Turbo Initech's library stops
% at the ADR-0007 DEC-05 boundary; DEC-05 itself names the integer-to-
% decimal conversion as hand-written in-subset Pascal.  These are user
% routines (not registered as standard routines), each a complete program
% replayed by test-pubs-turbo.  Idioms per ADR-0007 DEC-03 (nested-if guards,
% flag loops).
\Gates{test-compiler, test-compiler-os, test-pubs-turbo}
\chapter{Routines written in Pascal}
\ix{routines, written in Pascal}

The routines of this appendix are not part of Turbo Initech.  They are
the routines that other Pascal libraries provide and that Turbo Initech
programs write for themselves, each in a complete program you can compile
as it stands.  Copy the routine you need into your own program.  Each one
follows the rules of complete evaluation: every test that protects an
index is a separate \rw{if}.

\section{Converting an integer to a string}
\ix{integer to string}
\Program{inttostr}
\ProgramOutput{inttostr}

The digits are produced from the last to the first with \rw{mod} and
\rw{div} and are placed in front of the string built so far.  The routine
is correct for every integer except $-2147483648$, whose negation is not an
integer.

\section{Converting a string to an integer}
\ix{string to integer}
\Program{strtoint}
\ProgramOutput{strtoint}

\section{Converting letters to capitals}
\ix{capital letters}
\Program{upcase}
\ProgramOutput{upcase}

The nested \rw{if} statements test the two bounds separately; with \rw{and}
the routine would do the same, since both operands are characters, but
the nested form is the one that also works where the second test may only
be made when the first is true.

\section{Taking part of a string and finding a character}
\ix{substrings}
\Program{copy}
\ProgramOutput{copy}

\id{Copy} tests each index against both ends of the string before it reads
the character, so that a request past the end is safe.  \id{Pos} stops at
the first match with a Boolean flag.

\section{Reading lines from a file}
\ix{reading lines}
\ix{files, lines}
\Program{readline}
\ProgramOutput{readline}

\id{ReadLine} reads one byte at a time into a string of capacity 1 and ends
a line at the line-feed character.  Because a file cannot be a parameter,
the routine reads the global file variable \id{Data}.

\section{Inserting and deleting characters}
\ix{inserting characters}
\ix{deleting characters}
\Program{insert}
\ProgramOutput{insert}

Both routines build the new value in a local string and assign it back, so
that no character is overwritten before it has been copied.

\section{Removing blanks}
\ix{blanks, removing}
\Program{trim}
\ProgramOutput{trim}

\section{Comparing without regard to case}
\ix{comparing strings}
\Program{sametext}
\ProgramOutput{sametext}

The relational operators compare character codes, so \code{'Initech'} and
\code{'INITECH'} differ; \id{SameText} compares the letters after changing
them to small letters.

\section{Hexadecimal}
\ix{hexadecimal}
\Program{hex}
\ProgramOutput{hex}

Turbo Initech has no hexadecimal constants, but a program can write numbers
in hexadecimal, such as the load address of a program, 40100h.
